% =============================================================
% Pilot Validation
% =============================================================

\section{Pilot Validation}

A small pilot study was carried out during April--May 2026 to test whether
the labs work well in practice and whether students can complete them using
the provided guides.

\subsection{Pilot Design}

Five participants took part in the pilot. All of them are currently studying
cybersecurity, so they have some technical background, but they are not
experts in all the topics covered by the labs. The author was available
during the sessions to answer questions when participants got stuck, so the
study reflects a guided but realistic classroom setting rather than fully
independent self-study.

One participant has documented dyslexia. This made it possible to get some
early feedback on the accessibility features of the guides, such as the use
of OpenDyslexic font, the structured task layout, and the hint system.

Each participant recorded how long they took to complete their lab and
reported any technical problems they encountered. After finishing, they
filled in a short questionnaire with four questions rated on a 1--5 scale
(clarity of the guide, difficulty, overall experience, and whether the guide
was enough to complete the lab), plus a few open questions about what was
hardest and what they would change.

\subsection{Quantitative Results}

Table~\ref{tab:pilot_results} shows the average scores and times for each
lab.

\begin{table}[H]
  \centering
  \small
  \setlength{\tabcolsep}{5pt}
  \renewcommand{\arraystretch}{1.15}
  \begin{tabularx}{\textwidth}{l c c c c c}
    \toprule
    \textbf{Lab} & \textbf{\textit{n}} & \textbf{Clarity} &
    \textbf{Difficulty} & \textbf{Experience} &
    \textbf{Avg.\ time (min)} \\
    \midrule
    Lab~1 --- Reconnaissance      & 3 & 4.33 & 2.33 & 4.67 & 88  \\
    Lab~2 --- Web Vulnerabilities  & 3 & 5.00 & 3.00 & 4.67 & 95  \\
    Lab~3 --- Incident Response    & 3 & 4.00 & 3.67 & 4.67 & 122 \\
    \midrule
    \textbf{Overall}              & 9 & \textbf{4.44} & \textbf{3.00} &
                                       \textbf{4.67} & \textbf{102} \\
    \bottomrule
  \end{tabularx}
  \caption{Pilot study results --- mean scores (1--5 Likert) and resolution times}
  \label{tab:pilot_results}
\end{table}

All five participants completed their labs. No technical problem stopped
anyone from finishing. The average satisfaction score was \textbf{4.67 out
of 5}, which is above the 4/5 target set in OE4 (Section~3). The average
time across all labs was \textbf{102 minutes}, which is within the 90--120
minute target.

Lab~1 took an average of 88 minutes, just under the target. This is
expected since it is the first lab and the most guided one. Lab~3 took an
average of 122 minutes, just over the target. This makes sense because it
includes a forensic analysis task and a written incident report, which take
more time than a purely technical task. Both differences are small and fit
the planned progression from easier to harder labs.

\subsection{Qualitative Analysis}

The open questions showed some common patterns across the three labs.

\paragraph{Getting started in the environment.}
Two participants had trouble at the beginning before the technical tasks even
started. One spent around ten minutes trying to connect via SSH to the
attacker machine in Lab~1, not knowing it only has a VNC console inside
EVE-NG. Another participant in Lab~2 sent a payload to the wrong server
because the guide did not explain upfront that nginx works as a reverse proxy
in front of two backend servers. Both cases suggest that each student guide
needs a short section at the beginning explaining how to access the lab and
what the key parts of the topology are.

\paragraph{Understanding why, not just how.}
Several participants said they found it more helpful when the guide explained
the reason behind a step, not just the command to run. In Lab~2, the choice
of listener IP address was only explained in a late hint, so one participant
did not understand it until that point. In Lab~3, using \texttt{find} to
escalate privileges through GTFOBins\footnote{GTFOBins is a list of Unix
binaries that can be used to bypass local security restrictions:
\url{https://gtfobins.github.io}.} felt confusing without any context about
what GTFOBins is. Adding a short explanation in the task body --- not just
in hints --- would make the labs more educational.

\paragraph{Log readability.}
The participant with dyslexia found the \texttt{auth.log} output in Lab~3
very hard to read. Log lines with timestamps, hostnames, PIDs, and messages
all together are a lot of information at once. The guide assumed students
already knew how to use \texttt{grep} to filter the relevant lines. Adding
a short introduction to log structure and showing one example of expected
output for each \texttt{grep} command would help all students, not just
those with reading difficulties.

\paragraph{Connecting the three labs.}
The participant who completed all three labs pointed out that the overall
flow makes sense (reconnaissance, then web attack, then incident response)
but felt that the connection between the labs could be stronger. Framing
Lab~3 as the investigation of an incident caused by the same type of attack
seen in Labs~1 and~2 would help students see the full picture --- attack and
defence as two sides of the same chain.

\subsection{KPI Compliance}

\begin{table}[H]
  \centering
  \small
  \renewcommand{\arraystretch}{1.15}
  \begin{tabularx}{\textwidth}{X c c c}
    \toprule
    \textbf{KPI} & \textbf{Target} & \textbf{Result} & \textbf{Status} \\
    \midrule
    Mean satisfaction score      & $\geq 4/5$   & 4.67/5  & $\checkmark$ \\
    Lab completion rate          & 100\%         & 9/9     & $\checkmark$ \\
    Mean resolution time         & 90--120 min   & 102 min & $\checkmark$ \\
    Blocking technical incidents & 0             & 0       & $\checkmark$ \\
    \bottomrule
  \end{tabularx}
  \caption{OE4 KPI compliance after pilot study}
  \label{tab:kpi_compliance}
\end{table}

All four KPIs from OE4 were met. The small differences in per-lab times are
not a problem: the overall average is within target, and the pattern matches
the intended difficulty progression of the lab sequence.

\subsection{Identified Improvements}

The pilot showed a number of small issues. None of them stopped anyone from
finishing, but fixing them would make the labs clearer and more useful:

\begin{itemize}
  \item \textbf{All labs}: add a short \textit{Environment Access} section
        at the start of each student guide explaining console types (VNC
        vs.\ SSH) and the key parts of the topology.
  \item \textbf{Lab~2}: explain the listener IP choice in the task body
        instead of a hint; add a note about the nginx reverse proxy before
        the task sequence; remind students to check that the listener is
        active before posting the payload.
  \item \textbf{Lab~3}: add a note about boot time before the first SSH
        connection; clarify that the database server is only reachable from
        inside the web server; add one line of expected \texttt{grep} output
        per command; add a short description of GTFOBins; add a reflection
        question at the end of the final task about real-world detection.
\end{itemize}
